\chapter{Correções do parecer de orientação (2026-10-29)}
\label{ap:parecer}

% Fonte: research/next-steps/parecer-orientador-2026-10-29.md (vault).
% Código: branch feature/wrap-up-review; decisões 53--60 em experiments/DECISIONS.md.
% Estado: feito = no código ou no texto; pendente (servidor) = depende do serra1,
% indisponível desde 2026-10-21; pendente (orientador/autor) = decisão fora do código.

\section{Críticos}
\begin{itemize}
\item C1, procedência e ética dos dados.
  \begin{itemize}
  \item Feito: origem registrada como ``vídeos públicos na internet'' (Seção~\ref{cap:experimentos}); cópias locais de vídeos e resultados apagadas; tratamento de dados corrigido (Seção~\ref{cap:metodo}).
  \item Pendente (autor): fonte, licença e termos de cada vídeo.
  \item Pendente (orientador): consulta ao CEP; plano B com gravações próprias consentidas.
  \end{itemize}
\item C2, tamanho e representatividade dos dados.
  \begin{itemize}
  \item Feito: valores por vídeo, macro-média e múltiplo do acaso (Tabela~\ref{tab:derived-chance}); duração e participação de cada vídeo no agregado; ameaça à validade externa (Seção~\ref{sec:ameacas-validade}).
  \item Pendente: mais dados; segundo reconhecedor não usado na otimização.
  \end{itemize}
\item C3, hiperparâmetros escolhidos nos dados de avaliação.
  \begin{itemize}
  \item Feito: LOVO para $\beta$ (Tabela~\ref{tab:derived-lovo}); ameaça à validade interna declarada.
  \item Pendente: vídeos de teste separados com o pipeline congelado.
  \end{itemize}
\item C4, modelo de ameaça e alegação de anonimização.
  \begin{itemize}
  \item Feito: modelo de ameaça (Seção~\ref{sec:ameaca}); alegação trocada para desidentificação facial em resumo, abstract, palavras-chave, introdução, discussão e conclusão; identificadores fora do rosto nas limitações.
  \item Pendente: \emph{parrot attack} (trabalho futuro); título (orientador).
  \end{itemize}
\item C5, furo na garantia \emph{fail-closed}.
  \begin{itemize}
  \item Feito: fração mínima de pixels alterados (1\%) no Estágio~2 e na composição; teste do invariante em quadros sintéticos.
  \item Pendente (servidor): frequência do \emph{fail-open} nas execuções existentes; anotação de rostos nunca detectados (\texttt{miss\_sheet}).
  \end{itemize}
\item C6, validade da utilidade.
  \begin{itemize}
  \item Feito: proxy de tarefa (orientação, $\kappa$) implementado; ``sem custo'' reescrito como ``sem diferença significativa''; ameaça à validade de construto.
  \item Pendente (servidor): $\kappa$ de orientação. Pendente: TOST; avaliação humana com gravação consentida; título (``Expression-Preserving'').
  \end{itemize}
\end{itemize}

\section{Importantes}
\begin{itemize}
\item I1, visões de privacidade inconsistentes: divergência declarada (Seção~\ref{cap:experimentos}). Pendente (servidor): E2/E3 com \texttt{--video-view}.
\item I2, ocultar infla a privacidade: declarado; par (rank-1 nos mostrados, parcela mostrada) na discussão.
\item I3, rostos grandes e o professor: estratos por tamanho no E3e e no E5 implementados; pendente (servidor). Política por tamanho: trabalho futuro.
\item I4, identidade sintética nasce do rosto real: declarado no modelo de ameaça; E4-lite como trabalho futuro.
\item I5, estatística: Holm por família; tabelas por vídeo e macro. Pendente: valores ajustados (próxima exportação); TOST.
\item I6, escopo da métrica de custo: definições de quadros/s (E2) e ms por rosto (E5) separadas; projeção rotulada; exclusões listadas. Pendente: tempo ponta a ponta por estágio.
\item I7, reprodutibilidade: \texttt{SOURCES.md} com commit da ponte, Python, torch e \texttt{uv.lock}; E5 aborta sem o cache do E3. Pendente (servidor): push da ponte, E0/E1 no commit final, tag \texttt{tcc-defesa}.
\item I8, sem testes: 45 testes unitários (\texttt{tests/}).
\item I9, dados reais persistidos: crops de rostos não substituídos não gravados; scratch apagado em SIGTERM e SIGHUP; \texttt{purge.sh}.
\item I10, novidade modesta: contribuições enquadradas como empíricas e de sistema (Seção~\ref{cap:intro}).
\item I11, dez divergências entre texto e código, corrigidas no texto:
  \begin{enumerate}
  \item dilatação só em tracks confirmados;
  \item rastreador sem modelo de movimento, associação gulosa;
  \item margem de contexto 0,8 do BLANKET não calibrada;
  \item $\beta = 0{,}35$ não equivale ao mix nativo (nativo normalizado $\approx 0{,}26$);
  \item re-detecção a 320\,px e corte 0,5;
  \item E1 com limiar 0,3;
  \item cache de identidades chaveado pela seed derivada do caminho;
  \item apagamento do scratch e dos crops descrito como realmente ocorre;
  \item GPU-horas sem detecção, pré-passe, identidades e composição;
  \item quadros/s do E2 com outra definição (tempo total).
  \end{enumerate}
\end{itemize}

\section{Menores e formato}
\begin{itemize}
\item Objetivos geral e específicos (O1--O4, com situação); classificação da pesquisa; desvios e Quadro~\ref{qua:desvios}.
\item ``Cap.'' substituído por ``Seção''.
\item Listas de figuras, tabelas e quadros ativadas; todo float com título acima, \texttt{\textbackslash label}, fonte abaixo e menção no texto.
\item Tabelas em português com vírgula decimal (\texttt{generated/<experimento>/pt}). Pendente (servidor): eixos das figuras em português.
\item Resumo e abstract em 3ª pessoa, parágrafo único, sem rótulos internos; palavras-chave com ``;'' e minúsculas.
\item Siglas completadas (CEP, CPU, GPU, EMA, IoU, LOVO, RMS, TCC, TOST, VAE e outras); $\tau$ desambiguado em $\tau_h$ e $\tau_a$; $\alpha_f$ na lista de símbolos.
\item Bibliografia: entradas ``verificar'' conferidas (LFW, FDDB, Deep OC-SORT, MPGD, HSEmotion, Fantômas); ``Disponível em'' e ``Acesso em'' nas referências online.
\item Restos do modelo removidos: \texttt{a1-exemplo}, \texttt{a2-exemplo}, \texttt{exemplo.c}, \texttt{exemplo-figura.pdf}, \texttt{main\_maziero.pdf}, \texttt{lipsum}; ficha catalográfica fora (não exigida para TCC); descrição do trabalho pelo modelo SiBi; ``Prof. Dr.'' no orientador.
\item README atualizado; \texttt{phase3\_temporal} removido; \texttt{fps} do \texttt{tracks.json}; código morto removido.
\item Pendente (orientador): idioma do documento e título; metas da hipótese.
\item Pendente (autor): agradecimentos e dedicatória; desenvolver os tópicos em parágrafos; diagrama do pipeline; geometria do P2.
\end{itemize}
